\documentclass[12pt,a4paper]{article}\ref{}

% ---------- Packages ----------
\usepackage[margin=1in]{geometry}
\usepackage{times}
\usepackage{setspace}
\usepackage{enumitem}
\usepackage{array}
\usepackage{longtable}
\usepackage{booktabs}
\usepackage{tabularx}
\usepackage{ragged2e}
\usepackage{hyperref}
\usepackage{xcolor}
\usepackage{amsmath}
\usepackage{fontenc}
\usepackage[utf8]{inputenc}

% ---------- Formatting ----------
\setstretch{1.15}

\setlength{\parindent}{0pt}
\setlength{\parskip}{6pt}

\hypersetup{
    colorlinks=true,
    linkcolor=black,
    urlcolor=blue
}

% ---------- Custom commands ----------
\newcommand{\sectiontitle}[1]{%
    \vspace{8pt}
    \textbf{\large #1}
    \vspace{4pt}
}

\newcommand{\subsectiontitle}[1]{%
    \vspace{5pt}
    \textbf{#1}
    \vspace{3pt}
}

\begin{document}

% =========================================================
% TITLE PAGE / HEADER
% =========================================================

\begin{center}

{\Large \textbf{SmartExpense: Event-Aware Expense Tracker with AI-Powered Spend Prediction}}

\vspace{8pt}

{\large \textbf{Project Proposal — UCS503P}}

\vspace{12pt}

Submitted to: Mr. Asif

\vspace{4pt}

Thapar Institute of Engineering and Technology

\vspace{12pt}

Divanshi Garg [1024030441], Navya Arora [1024030438]

\vspace{4pt}

August 2026

\end{center}

\vspace{12pt}

% =========================================================
% 1. INTRODUCTION
% =========================================================

\sectiontitle{1. Introduction and Problem Statement}

Personal finance tools are mostly passive, ie they record past expenses and show
historical charts but provide little guidance for future spending. Predictable events
such as birthdays, anniversaries, and festivals can still cause unexpected financial
pressure because existing tools rarely connect upcoming events with spending
patterns.

The main gaps are:

\begin{itemize}[leftmargin=1.5em]
    \item No event-aware forecasting of future expenses.
    \item Fragmented planning, with budgets, notes and reminders handled in different apps.
    \item Manual expense entry, which can be time-consuming and inconvenient.
\end{itemize}

SmartExpense addresses these gaps through an integrated platform that combines
expense tracking, event-aware spending prediction, reminders, and a Notes-style
planning workspace. It aims to transform expense tracking from a passive
record-keeping tool into a proactive financial planning system.

% =========================================================
% 2. OBJECTIVES
% =========================================================

\sectiontitle{2. Objectives}

\textbf{Primary Goal:} To design, build, and deploy a full-stack web application that predicts
future-month personal expenditure per category, adjusted for user-registered life
events, and notifies users proactively before those events occur.

\begin{itemize}[leftmargin=1.5em]
    \item To implement a secure MERN-stack expense tracking system with JWT-based
    authentication, transaction CRUD, and keyword-based categorization supporting
    both manual entry and bulk CSV upload.

    \item To build a Python/FastAPI microservice that fits a Prophet time-series model
    on per-user transaction history, with registered events as binary regressors, and
    exposes a \texttt{/predict-spend} endpoint called monthly by a scheduled job.

    \item To implement an event registry module that stores recurring and one-time life
    events, links them to spending categories, and triggers proactive email reminders
    (Nodemailer) within 7 days of each event.

    \item To deliver an integrated Notes and Saved Items module where users can plan
    for events.

    \item To deploy all three services (React frontend, Express backend, FastAPI ML
    service) to production on free-tier cloud infrastructure (Vercel, Render,
    MongoDB Atlas) with a live accessible URL.
\end{itemize}

% =========================================================
% 3. METHODOLOGY
% =========================================================

\sectiontitle{3. Methodology}

System architecture is defined upfront, features are built in three sequential
iterations with deployment after each, and the ML layer is validated against
measurable accuracy targets before being integrated into the product.

\subsectiontitle{3.1 System Architecture and Tech Stack}

The application is a three-service architecture:

\vspace{5pt}

\begin{center}
\renewcommand{\arraystretch}{1.3}
\begin{tabularx}{\textwidth}{
    >{\raggedright\arraybackslash}p{0.18\textwidth}
    >{\raggedright\arraybackslash}p{0.42\textwidth}
    >{\raggedright\arraybackslash}X
}
\toprule
\textbf{Layer} & \textbf{Technology} & \textbf{Role} \\
\midrule

Frontend &
React (Vite) + Tailwind CSS + Recharts &
UI, routing, data visualization \\

Backend API &
Node.js + Express + JWT + bcrypt + node-cron &
Auth, CRUD, scheduling, orchestration \\

ML Service &
Python + FastAPI + Prophet + scikit-learn &
Spend forecasting, anomaly detection \\

Database &
MongoDB Atlas + Mongoose &
All persistent storage \\

\bottomrule
\end{tabularx}
\end{center}

\subsectiontitle{3.2 Module Development}

\textbf{Auth Module:} bcrypt password hashing, JWT-based stateless auth, protected routes
on all data endpoints.

\textbf{Transactions Module:} Full CRUD for transactions with amount, date, merchant, and
category fields. CSV bulk upload (Multer) parses bank statement exports; a
keyword-matching function (e.g., ``swiggy/zomato'' : Food, ``uber/ola'' : Transport)
assigns categories automatically, covering 80\% of transactions without manual tagging.

\textbf{Events Module:} Users register recurring or one-time events (name, date, recurrence,
linked category, optional expected spend). MongoDB stores event history including
actual post-event spend for model learning.

\textbf{AI/Prediction Service (FastAPI):} On a monthly cron trigger, Express calls
POST /retrain-and-predict with the user's transaction history. FastAPI aggregates
transactions to monthly category totals, adds binary event regressors for the upcoming
month, and fits a Prophet model. The resulting spend prediction per category is stored
in a Predictions collection in MongoDB. A z-score threshold on individual transactions
flags anomalies for display in the dashboard.

\textit{// `cron' here just means time}

\textbf{Notes and Saved Items Module:} Freeform notes (text, checkboxes, pinning) linked to
events. Saved item costs are summed and passed as an additional input to the prediction
service. (if users click on the option of adding to the spendings)

\subsectiontitle{3.3 Data Collection and Model Training}

Training data is sourced from the user's own transaction history, entered via manual
logging or CSV upload from their bank or UPI app statement. For the initial development
and pilot phase, the developer's own anonymized 3-month bank statement will seed the
system. A cold-start fallback (category averages from the first 1--2 months) handles new
users before sufficient history exists to fit a reliable model. Monthly cron-triggered
per-user retraining ensures the model improves as transaction history grows.

\subsectiontitle{3.4 Testing and Validation}

Each module will be validated with Postman API tests before integration. End-to-end
testing will cover the critical path: CSV upload, categorization, event registration,
prediction trigger, reminder email delivery. The primary ML accuracy metric is Mean
Absolute Error (MAE) between predicted and actual monthly category spend, measured
at month-end from database records. Target: median MAE $\leq$ ₹500 per category after
3 months of data.

\subsectiontitle{3.5 Risks and Mitigations}

\begin{itemize}[leftmargin=1.5em]
    \item \textbf{Cold-start (no history for new users):} category-average fallback,
    communicated clearly in the UI; CSV import at onboarding accelerates history
    collection.

    \item \textbf{Manual entry abandonment:} weekly ``you haven't logged in X days''
    nudge via the existing reminder infrastructure; CSV upload as a frictionless
    alternative.

    \item \textbf{Render free-tier concurrent retraining load:} stagger monthly cron
    jobs across users rather than triggering all at midnight on the 1st.
\end{itemize}

% =========================================================
% 4. TIMELINE
% =========================================================
\vspace{5pt}


\newpage
\sectiontitle{4. Timeline}

\vspace{5pt}

\begin{center}
\renewcommand{\arraystretch}{1.35}

\begin{longtable}{
    >{\raggedright\arraybackslash}p{0.09\textwidth}
    >{\raggedright\arraybackslash}p{0.17\textwidth}
    >{\raggedright\arraybackslash}p{0.39\textwidth}
    >{\raggedright\arraybackslash}p{0.27\textwidth}
}

\toprule
\textbf{Week} &
\textbf{Phase} &
\textbf{Tasks} &
\textbf{Deliverable} \\
\midrule
\endfirsthead

\toprule
\textbf{Week} &
\textbf{Phase} &
\textbf{Tasks} &
\textbf{Deliverable} \\
\midrule
\endhead

3--4 &
Setup \& Planning &
Finalize requirements, set up GitHub repository, create MERN project structure,
connect MongoDB, and design database schema &
Working project setup and database connection \\

\midrule

5--6 &
Expense Module &
Implement expense CRUD, categories, transaction forms, and basic dashboard/charts &
Functional expense tracking module \\

\midrule

7--8 &
Events \& Reminders &
Implement event CRUD, recurring events, and scheduled reminders using node-cron &
Functional event and reminder system \\

\midrule

9--10 &
AI/ML Module &
Set up Python/FastAPI service, prepare transaction data, develop initial forecasting
model, and integrate it with the backend &
Initial spending prediction integrated with dashboard \\

\midrule

11 &
Integration \& Testing &
Integrate all modules, test APIs, fix bugs, improve UI, and evaluate prediction results &
Integrated and tested application \\

\midrule

12 &
Finalization \& Deployment &
Final testing, documentation, deployment, and presentation preparation &
Final working application and project documentation \\

\bottomrule
\end{longtable}
\end{center}

% =========================================================
% 5. RESOURCES AND BUDGET
% =========================================================

\sectiontitle{5. Resources and Budget}

\textbf{Human Resources:} Group of two project.

\begin{itemize}[leftmargin=1.5em]
    \item Mern stack application and ideation is handled by Divanshi Garg
    \item ML Prediction of future spendings and ideation handled by Navya Arora
\end{itemize}

\textbf{Software and Tools:} All tools are free and open-source — React, Node.js,
Express, Python, FastAPI, Prophet, scikit-learn, MongoDB Atlas (free tier), Vercel
(free tier), Render (free tier), Postman (free), GitHub.

\textbf{External APIs:} Gmail SMTP via Nodemailer for reminder emails (free).

\textbf{Budget:} No external funding required. All hosting, database, and API services
operate within free tiers appropriate for a portfolio-scale educational deployment.
No hardware beyond a personal laptop is needed.

% =========================================================
% 6. SUMMARY
% =========================================================

\sectiontitle{6. Summary}

\textbf{The Gap:} Existing expense trackers are passive, they record spending but do not
connect it to a user's calendar of predictable life events, leaving users financially
unprepared for costs they could have anticipated.

\textbf{The Solution:} SmartExpense is a three-service MERN + Python web application
that forecasts next-month spending per category using Predictive time-series models
with event-based regressors, sends proactive reminders ahead of registered events,
and provides an integrated notes workspace for event planning.

\textbf{The Impact:} Users shift from reactive expense logging to proactive financial
planning: knowing before a birthday month arrives what to budget, where to shop,
and how much they typically spend.

\end{document}
